% Options for packages loaded elsewhere
\PassOptionsToPackage{unicode}{hyperref}
\PassOptionsToPackage{hyphens}{url}
\documentclass[
]{article}
\usepackage{xcolor}
\usepackage[margin=1in]{geometry}
\usepackage{amsmath,amssymb}
\setcounter{secnumdepth}{-\maxdimen} % remove section numbering
\usepackage{iftex}
\ifPDFTeX
  \usepackage[T1]{fontenc}
  \usepackage[utf8]{inputenc}
  \usepackage{textcomp} % provide euro and other symbols
\else % if luatex or xetex
  \usepackage{unicode-math} % this also loads fontspec
  \defaultfontfeatures{Scale=MatchLowercase}
  \defaultfontfeatures[\rmfamily]{Ligatures=TeX,Scale=1}
\fi
\usepackage{lmodern}
\ifPDFTeX\else
  % xetex/luatex font selection
\fi
% Use upquote if available, for straight quotes in verbatim environments
\IfFileExists{upquote.sty}{\usepackage{upquote}}{}
\IfFileExists{microtype.sty}{% use microtype if available
  \usepackage[]{microtype}
  \UseMicrotypeSet[protrusion]{basicmath} % disable protrusion for tt fonts
}{}
\makeatletter
\@ifundefined{KOMAClassName}{% if non-KOMA class
  \IfFileExists{parskip.sty}{%
    \usepackage{parskip}
  }{% else
    \setlength{\parindent}{0pt}
    \setlength{\parskip}{6pt plus 2pt minus 1pt}}
}{% if KOMA class
  \KOMAoptions{parskip=half}}
\makeatother
\usepackage{color}
\usepackage{fancyvrb}
\newcommand{\VerbBar}{|}
\newcommand{\VERB}{\Verb[commandchars=\\\{\}]}
\DefineVerbatimEnvironment{Highlighting}{Verbatim}{commandchars=\\\{\}}
% Add ',fontsize=\small' for more characters per line
\usepackage{framed}
\definecolor{shadecolor}{RGB}{248,248,248}
\newenvironment{Shaded}{\begin{snugshade}}{\end{snugshade}}
\newcommand{\AlertTok}[1]{\textcolor[rgb]{0.94,0.16,0.16}{#1}}
\newcommand{\AnnotationTok}[1]{\textcolor[rgb]{0.56,0.35,0.01}{\textbf{\textit{#1}}}}
\newcommand{\AttributeTok}[1]{\textcolor[rgb]{0.13,0.29,0.53}{#1}}
\newcommand{\BaseNTok}[1]{\textcolor[rgb]{0.00,0.00,0.81}{#1}}
\newcommand{\BuiltInTok}[1]{#1}
\newcommand{\CharTok}[1]{\textcolor[rgb]{0.31,0.60,0.02}{#1}}
\newcommand{\CommentTok}[1]{\textcolor[rgb]{0.56,0.35,0.01}{\textit{#1}}}
\newcommand{\CommentVarTok}[1]{\textcolor[rgb]{0.56,0.35,0.01}{\textbf{\textit{#1}}}}
\newcommand{\ConstantTok}[1]{\textcolor[rgb]{0.56,0.35,0.01}{#1}}
\newcommand{\ControlFlowTok}[1]{\textcolor[rgb]{0.13,0.29,0.53}{\textbf{#1}}}
\newcommand{\DataTypeTok}[1]{\textcolor[rgb]{0.13,0.29,0.53}{#1}}
\newcommand{\DecValTok}[1]{\textcolor[rgb]{0.00,0.00,0.81}{#1}}
\newcommand{\DocumentationTok}[1]{\textcolor[rgb]{0.56,0.35,0.01}{\textbf{\textit{#1}}}}
\newcommand{\ErrorTok}[1]{\textcolor[rgb]{0.64,0.00,0.00}{\textbf{#1}}}
\newcommand{\ExtensionTok}[1]{#1}
\newcommand{\FloatTok}[1]{\textcolor[rgb]{0.00,0.00,0.81}{#1}}
\newcommand{\FunctionTok}[1]{\textcolor[rgb]{0.13,0.29,0.53}{\textbf{#1}}}
\newcommand{\ImportTok}[1]{#1}
\newcommand{\InformationTok}[1]{\textcolor[rgb]{0.56,0.35,0.01}{\textbf{\textit{#1}}}}
\newcommand{\KeywordTok}[1]{\textcolor[rgb]{0.13,0.29,0.53}{\textbf{#1}}}
\newcommand{\NormalTok}[1]{#1}
\newcommand{\OperatorTok}[1]{\textcolor[rgb]{0.81,0.36,0.00}{\textbf{#1}}}
\newcommand{\OtherTok}[1]{\textcolor[rgb]{0.56,0.35,0.01}{#1}}
\newcommand{\PreprocessorTok}[1]{\textcolor[rgb]{0.56,0.35,0.01}{\textit{#1}}}
\newcommand{\RegionMarkerTok}[1]{#1}
\newcommand{\SpecialCharTok}[1]{\textcolor[rgb]{0.81,0.36,0.00}{\textbf{#1}}}
\newcommand{\SpecialStringTok}[1]{\textcolor[rgb]{0.31,0.60,0.02}{#1}}
\newcommand{\StringTok}[1]{\textcolor[rgb]{0.31,0.60,0.02}{#1}}
\newcommand{\VariableTok}[1]{\textcolor[rgb]{0.00,0.00,0.00}{#1}}
\newcommand{\VerbatimStringTok}[1]{\textcolor[rgb]{0.31,0.60,0.02}{#1}}
\newcommand{\WarningTok}[1]{\textcolor[rgb]{0.56,0.35,0.01}{\textbf{\textit{#1}}}}
\usepackage{graphicx}
\makeatletter
\newsavebox\pandoc@box
\newcommand*\pandocbounded[1]{% scales image to fit in text height/width
  \sbox\pandoc@box{#1}%
  \Gscale@div\@tempa{\textheight}{\dimexpr\ht\pandoc@box+\dp\pandoc@box\relax}%
  \Gscale@div\@tempb{\linewidth}{\wd\pandoc@box}%
  \ifdim\@tempb\p@<\@tempa\p@\let\@tempa\@tempb\fi% select the smaller of both
  \ifdim\@tempa\p@<\p@\scalebox{\@tempa}{\usebox\pandoc@box}%
  \else\usebox{\pandoc@box}%
  \fi%
}
% Set default figure placement to htbp
\def\fps@figure{htbp}
\makeatother
\setlength{\emergencystretch}{3em} % prevent overfull lines
\providecommand{\tightlist}{%
  \setlength{\itemsep}{0pt}\setlength{\parskip}{0pt}}
\usepackage{bookmark}
\IfFileExists{xurl.sty}{\usepackage{xurl}}{} % add URL line breaks if available
\urlstyle{same}
\hypersetup{
  pdftitle={Relatório Técnico: Cenários de Estimativas de Emissões de GEE no Brasil com o Modelo MAgPIE},
  pdfauthor={Fernanda Valente, Sabrina de Matos Carlos, Cícero Lima, Ângelo Gurgel},
  hidelinks,
  pdfcreator={LaTeX via pandoc}}

\title{Relatório Técnico: Cenários de Estimativas de Emissões de GEE no
Brasil com o Modelo MAgPIE}
\author{Fernanda Valente, Sabrina de Matos Carlos, Cícero Lima, Ângelo
Gurgel}
\date{2026-03-03}

\begin{document}
\maketitle

\begin{Shaded}
\begin{Highlighting}[]
\CommentTok{\# Opções globais dos chunks}
\NormalTok{knitr}\SpecialCharTok{::}\NormalTok{opts\_chunk}\SpecialCharTok{$}\FunctionTok{set}\NormalTok{(}\AttributeTok{warning =} \ConstantTok{FALSE}\NormalTok{, }\AttributeTok{message =} \ConstantTok{FALSE}\NormalTok{, }\AttributeTok{fig.width =} \DecValTok{12}\NormalTok{, }\AttributeTok{fig.height =} \DecValTok{6}\NormalTok{, }\AttributeTok{out.width =} \StringTok{"100\%"}\NormalTok{)}

\CommentTok{\# Pacotes}
\FunctionTok{library}\NormalTok{(sf); }\FunctionTok{library}\NormalTok{(sp); }\FunctionTok{library}\NormalTok{(dplyr); }\FunctionTok{library}\NormalTok{(raster)}
\FunctionTok{library}\NormalTok{(madrat); }\FunctionTok{library}\NormalTok{(magclass); }\FunctionTok{library}\NormalTok{(magpie4); }\FunctionTok{library}\NormalTok{(luplot); }\FunctionTok{library}\NormalTok{(gdx2)}
\FunctionTok{library}\NormalTok{(terra); }\FunctionTok{library}\NormalTok{(ggplot2); }\FunctionTok{library}\NormalTok{(scales); }\FunctionTok{library}\NormalTok{(geobr); }\FunctionTok{library}\NormalTok{(rmdformats)}

\CommentTok{\# Paleta padrão para mapas}
\NormalTok{pal\_magpie }\OtherTok{\textless{}{-}} \FunctionTok{hcl.colors}\NormalTok{(}\DecValTok{20}\NormalTok{, }\StringTok{"YlOrBr"}\NormalTok{, }\AttributeTok{rev =} \ConstantTok{TRUE}\NormalTok{)}
\end{Highlighting}
\end{Shaded}

\begin{Shaded}
\begin{Highlighting}[]
\CommentTok{\# Mapa de clusters e filtro BRA com lon/lat numéricos}
\NormalTok{map\_file }\OtherTok{\textless{}{-}} \FunctionTok{readRDS}\NormalTok{(}\StringTok{"clustermap\_rev4.121\_c200\_67420\_h12.rds"}\NormalTok{)}

\NormalTok{BRA }\OtherTok{\textless{}{-}} \FunctionTok{subset}\NormalTok{(map\_file, country }\SpecialCharTok{==} \StringTok{"BRA"}\NormalTok{)}
\NormalTok{coords }\OtherTok{\textless{}{-}} \FunctionTok{strsplit}\NormalTok{(BRA}\SpecialCharTok{$}\NormalTok{cell, }\StringTok{"}\SpecialCharTok{\textbackslash{}\textbackslash{}}\StringTok{."}\NormalTok{)}
\NormalTok{mat }\OtherTok{\textless{}{-}} \FunctionTok{do.call}\NormalTok{(rbind, coords)}
\NormalTok{BRA}\SpecialCharTok{$}\NormalTok{lon }\OtherTok{\textless{}{-}} \FunctionTok{as.numeric}\NormalTok{(}\FunctionTok{gsub}\NormalTok{(}\StringTok{"p"}\NormalTok{, }\StringTok{"."}\NormalTok{, mat[,}\DecValTok{1}\NormalTok{]))}
\NormalTok{BRA}\SpecialCharTok{$}\NormalTok{lat }\OtherTok{\textless{}{-}} \FunctionTok{as.numeric}\NormalTok{(}\FunctionTok{gsub}\NormalTok{(}\StringTok{"p"}\NormalTok{, }\StringTok{"."}\NormalTok{, mat[,}\DecValTok{2}\NormalTok{]))}
\NormalTok{BRA}\SpecialCharTok{$}\NormalTok{iso }\OtherTok{\textless{}{-}}\NormalTok{ mat[,}\DecValTok{3}\NormalTok{]}

\CommentTok{\# Shapes}
\NormalTok{bra    }\OtherTok{\textless{}{-}}\NormalTok{ geobr}\SpecialCharTok{::}\FunctionTok{read\_country}\NormalTok{(}\AttributeTok{showProgress =} \ConstantTok{FALSE}\NormalTok{) }\SpecialCharTok{|\textgreater{}} \FunctionTok{st\_transform}\NormalTok{(}\DecValTok{4326}\NormalTok{)}
\NormalTok{biomes }\OtherTok{\textless{}{-}}\NormalTok{ geobr}\SpecialCharTok{::}\FunctionTok{read\_biomes}\NormalTok{(}\AttributeTok{showProgress =} \ConstantTok{FALSE}\NormalTok{)  }\SpecialCharTok{|\textgreater{}} \FunctionTok{st\_transform}\NormalTok{(}\DecValTok{4326}\NormalTok{)}
\NormalTok{state }\OtherTok{\textless{}{-}}\NormalTok{ geobr}\SpecialCharTok{::}\FunctionTok{read\_state}\NormalTok{(}\AttributeTok{showProgress =} \ConstantTok{FALSE}\NormalTok{) }\SpecialCharTok{|\textgreater{}} \FunctionTok{st\_transform}\NormalTok{(}\DecValTok{4326}\NormalTok{)}
\end{Highlighting}
\end{Shaded}

\begin{Shaded}
\begin{Highlighting}[]
\CommentTok{\# Formata "2010" {-}\textgreater{} "y2010.y2010"}
\NormalTok{fmt\_year }\OtherTok{\textless{}{-}} \ControlFlowTok{function}\NormalTok{(y) }\FunctionTok{paste0}\NormalTok{(}\StringTok{"y"}\NormalTok{, y, }\StringTok{".y"}\NormalTok{, y)}

\CommentTok{\# Converte magclass {-}\textgreater{} data.frame com nomes consistentes}
\NormalTok{mag\_to\_df }\OtherTok{\textless{}{-}} \ControlFlowTok{function}\NormalTok{(obj) \{}
\NormalTok{  df }\OtherTok{\textless{}{-}} \FunctionTok{as.data.frame}\NormalTok{(obj, }\AttributeTok{rev =} \DecValTok{2}\NormalTok{)}
  \FunctionTok{names}\NormalTok{(df) }\OtherTok{\textless{}{-}} \FunctionTok{c}\NormalTok{(}\StringTok{"lon"}\NormalTok{, }\StringTok{"lat"}\NormalTok{, }\StringTok{"iso"}\NormalTok{, }\StringTok{"year.start"}\NormalTok{, }\StringTok{"year.end"}\NormalTok{, }\StringTok{"emis\_source"}\NormalTok{, }\StringTok{"pollutants"}\NormalTok{, }\StringTok{"value"}\NormalTok{)}
\NormalTok{  df}
\NormalTok{\}}

\CommentTok{\# GWP (SAR 100y) padronizado}
\NormalTok{gwp\_sar }\OtherTok{\textless{}{-}} \FunctionTok{c}\NormalTok{(}\AttributeTok{co2 =} \DecValTok{1}\NormalTok{, }\AttributeTok{ch4 =} \DecValTok{21}\NormalTok{, }\AttributeTok{n2o =} \DecValTok{310}\NormalTok{)}

\NormalTok{mt\_to\_Gg }\OtherTok{\textless{}{-}} \ControlFlowTok{function}\NormalTok{(x) x }\SpecialCharTok{*} \DecValTok{1000}
\NormalTok{mt\_to\_kt }\OtherTok{\textless{}{-}} \ControlFlowTok{function}\NormalTok{(x) x }\SpecialCharTok{*} \DecValTok{1000}

\CommentTok{\# Aplica GWP e converte para Gg CO2e/ano}
\NormalTok{apply\_gwp }\OtherTok{\textless{}{-}} \ControlFlowTok{function}\NormalTok{(df, }\AttributeTok{gas =} \StringTok{"ch4"}\NormalTok{, }\AttributeTok{gwp\_table =}\NormalTok{ gwp\_sar) \{}
\NormalTok{  g }\OtherTok{\textless{}{-}} \FunctionTok{tolower}\NormalTok{(gas)}
  \ControlFlowTok{if}\NormalTok{ (}\SpecialCharTok{!}\NormalTok{g }\SpecialCharTok{\%in\%} \FunctionTok{names}\NormalTok{(gwp\_table)) }\FunctionTok{stop}\NormalTok{(}\StringTok{"Gás não encontrado na tabela de GWP."}\NormalTok{)}
\NormalTok{  df }\SpecialCharTok{|\textgreater{}}
    \FunctionTok{mutate}\NormalTok{(}\AttributeTok{value\_co2e =}\NormalTok{ value }\SpecialCharTok{*}\NormalTok{ gwp\_table[[g]])}
\NormalTok{\}}

\CommentTok{\# Rasteriza um data.frame (lon,lat,value) e retorna grid por ano (lista empilhada)}
\NormalTok{gridify\_by\_year }\OtherTok{\textless{}{-}} \ControlFlowTok{function}\NormalTok{(d, bra\_sf) \{}
  \FunctionTok{split}\NormalTok{(d, d}\SpecialCharTok{$}\NormalTok{year.end) }\SpecialCharTok{|\textgreater{}}
    \FunctionTok{lapply}\NormalTok{(}\ControlFlowTok{function}\NormalTok{(x) \{}
\NormalTok{      r  }\OtherTok{\textless{}{-}}\NormalTok{ terra}\SpecialCharTok{::}\FunctionTok{rast}\NormalTok{(x[, }\FunctionTok{c}\NormalTok{(}\StringTok{"lon"}\NormalTok{,}\StringTok{"lat"}\NormalTok{,}\StringTok{"value"}\NormalTok{)], }\AttributeTok{type =} \StringTok{"xyz"}\NormalTok{, }\AttributeTok{crs =} \StringTok{"EPSG:4326"}\NormalTok{)}
\NormalTok{      r  }\OtherTok{\textless{}{-}}\NormalTok{ terra}\SpecialCharTok{::}\FunctionTok{crop}\NormalTok{(r, terra}\SpecialCharTok{::}\FunctionTok{vect}\NormalTok{(bra\_sf))}
\NormalTok{      r  }\OtherTok{\textless{}{-}}\NormalTok{ terra}\SpecialCharTok{::}\FunctionTok{mask}\NormalTok{(r, terra}\SpecialCharTok{::}\FunctionTok{vect}\NormalTok{(bra\_sf))}
\NormalTok{      gd }\OtherTok{\textless{}{-}} \FunctionTok{as.data.frame}\NormalTok{(r, }\AttributeTok{xy =} \ConstantTok{TRUE}\NormalTok{, }\AttributeTok{na.rm =} \ConstantTok{TRUE}\NormalTok{)}
      \FunctionTok{names}\NormalTok{(gd) }\OtherTok{\textless{}{-}} \FunctionTok{c}\NormalTok{(}\StringTok{"lon"}\NormalTok{,}\StringTok{"lat"}\NormalTok{,}\StringTok{"value"}\NormalTok{)}
\NormalTok{      gd}\SpecialCharTok{$}\NormalTok{year.end }\OtherTok{\textless{}{-}} \FunctionTok{unique}\NormalTok{(x}\SpecialCharTok{$}\NormalTok{year.end)}
\NormalTok{      gd}
\NormalTok{    \}) }\SpecialCharTok{|\textgreater{}}
\NormalTok{    dplyr}\SpecialCharTok{::}\FunctionTok{bind\_rows}\NormalTok{()}
\NormalTok{\}}

\CommentTok{\# Soma total por fonte em um recorte (country, yearS) e retorna top{-}k nomes (dim k do magclass)}
\NormalTok{top\_k\_sources }\OtherTok{\textless{}{-}} \ControlFlowTok{function}\NormalTok{(dis\_obj, }\AttributeTok{country =} \StringTok{"BRA"}\NormalTok{, yearS, }\AttributeTok{k =} \DecValTok{2}\NormalTok{, }\AttributeTok{pattern =} \ConstantTok{NULL}\NormalTok{) \{}
\NormalTok{  ks }\OtherTok{\textless{}{-}}\NormalTok{ magclass}\SpecialCharTok{::}\FunctionTok{getItems}\NormalTok{(dis\_obj, }\AttributeTok{dim =} \DecValTok{3}\NormalTok{)}
  \ControlFlowTok{if}\NormalTok{ (}\SpecialCharTok{!}\FunctionTok{is.null}\NormalTok{(pattern)) ks }\OtherTok{\textless{}{-}} \FunctionTok{grep}\NormalTok{(pattern, ks, }\AttributeTok{value =} \ConstantTok{TRUE}\NormalTok{)}
\NormalTok{  totals }\OtherTok{\textless{}{-}} \FunctionTok{sapply}\NormalTok{(ks, }\ControlFlowTok{function}\NormalTok{(p) }\FunctionTok{sum}\NormalTok{(}\FunctionTok{as.numeric}\NormalTok{(dis\_obj[country, yearS, p]), }\AttributeTok{na.rm =} \ConstantTok{TRUE}\NormalTok{))}
\NormalTok{  totals }\OtherTok{\textless{}{-}}\NormalTok{ totals[totals }\SpecialCharTok{\textgreater{}} \DecValTok{0}\NormalTok{]}
  \FunctionTok{names}\NormalTok{(}\FunctionTok{sort}\NormalTok{(totals, }\AttributeTok{decreasing =} \ConstantTok{TRUE}\NormalTok{))[}\FunctionTok{seq\_len}\NormalTok{(}\FunctionTok{min}\NormalTok{(k, }\FunctionTok{length}\NormalTok{(totals)))]}
\NormalTok{\}}

\CommentTok{\# Extrai uma fonte específica e devolve df com lon/lat/iso/anos/value (Mt gás), mais \textquotesingle{}source\textquotesingle{}}
\NormalTok{extract\_one\_source }\OtherTok{\textless{}{-}} \ControlFlowTok{function}\NormalTok{(dis\_obj, country, yearS, src) \{}
\NormalTok{  obj }\OtherTok{\textless{}{-}}\NormalTok{ dis\_obj[country, yearS, src]}
\NormalTok{  d   }\OtherTok{\textless{}{-}} \FunctionTok{mag\_to\_df}\NormalTok{(obj)}
\NormalTok{  d}\SpecialCharTok{$}\NormalTok{lon   }\OtherTok{\textless{}{-}} \FunctionTok{as.numeric}\NormalTok{(}\FunctionTok{gsub}\NormalTok{(}\StringTok{"p"}\NormalTok{, }\StringTok{"."}\NormalTok{, d}\SpecialCharTok{$}\NormalTok{lon))}
\NormalTok{  d}\SpecialCharTok{$}\NormalTok{lat   }\OtherTok{\textless{}{-}} \FunctionTok{as.numeric}\NormalTok{(}\FunctionTok{gsub}\NormalTok{(}\StringTok{"p"}\NormalTok{, }\StringTok{"."}\NormalTok{, d}\SpecialCharTok{$}\NormalTok{lat))}
\NormalTok{  d}\SpecialCharTok{$}\NormalTok{source }\OtherTok{\textless{}{-}}\NormalTok{ src}
\NormalTok{  d}
\NormalTok{\}}

\CommentTok{\# Figura de mapa facetado por ano}
\NormalTok{plot\_faceted\_map }\OtherTok{\textless{}{-}} \ControlFlowTok{function}\NormalTok{(gdf, bra\_sf, biomes\_sf, title, }\AttributeTok{fill\_lab =} \StringTok{"kt CO2e"}\NormalTok{, }\AttributeTok{pal =}\NormalTok{ pal\_magpie) \{}
  \FunctionTok{ggplot}\NormalTok{() }\SpecialCharTok{+}
    \FunctionTok{geom\_raster}\NormalTok{(}\AttributeTok{data =}\NormalTok{ gdf, }\FunctionTok{aes}\NormalTok{(lon, lat, }\AttributeTok{fill =}\NormalTok{ value)) }\SpecialCharTok{+}
    \FunctionTok{geom\_sf}\NormalTok{(}\AttributeTok{data =}\NormalTok{ biomes\_sf, }\AttributeTok{fill =} \ConstantTok{NA}\NormalTok{, }\AttributeTok{color =} \StringTok{"grey35"}\NormalTok{, }\AttributeTok{linewidth =} \FloatTok{0.25}\NormalTok{) }\SpecialCharTok{+}
    \FunctionTok{geom\_sf}\NormalTok{(}\AttributeTok{data =}\NormalTok{ bra\_sf,    }\AttributeTok{fill =} \ConstantTok{NA}\NormalTok{, }\AttributeTok{color =} \StringTok{"grey15"}\NormalTok{, }\AttributeTok{linewidth =} \FloatTok{0.4}\NormalTok{) }\SpecialCharTok{+}
    \FunctionTok{coord\_sf}\NormalTok{(}\AttributeTok{xlim =} \FunctionTok{c}\NormalTok{(}\SpecialCharTok{{-}}\DecValTok{74}\NormalTok{, }\SpecialCharTok{{-}}\DecValTok{34}\NormalTok{), }\AttributeTok{ylim =} \FunctionTok{c}\NormalTok{(}\SpecialCharTok{{-}}\DecValTok{34}\NormalTok{, }\DecValTok{6}\NormalTok{), }\AttributeTok{expand =} \ConstantTok{FALSE}\NormalTok{) }\SpecialCharTok{+}
    \FunctionTok{facet\_wrap}\NormalTok{(}\SpecialCharTok{\textasciitilde{}}\NormalTok{ year.end) }\SpecialCharTok{+}
    \FunctionTok{scale\_fill\_gradientn}\NormalTok{(}\AttributeTok{colors =}\NormalTok{ pal, }\AttributeTok{name =}\NormalTok{ fill\_lab) }\SpecialCharTok{+}
    \FunctionTok{labs}\NormalTok{(}\AttributeTok{title =}\NormalTok{ title, }\AttributeTok{caption =} \StringTok{"Source: MAgPIE / Default version (SSP2)"}\NormalTok{) }\SpecialCharTok{+}
    \FunctionTok{theme\_minimal}\NormalTok{(}\AttributeTok{base\_size =} \DecValTok{12}\NormalTok{) }\SpecialCharTok{+}
    \FunctionTok{theme}\NormalTok{(}
      \AttributeTok{legend.position =} \StringTok{"right"}\NormalTok{,}
      \AttributeTok{panel.grid      =} \FunctionTok{element\_blank}\NormalTok{(),}
      \AttributeTok{plot.title      =} \FunctionTok{element\_text}\NormalTok{(}\AttributeTok{face =} \StringTok{"bold"}\NormalTok{, }\AttributeTok{size =} \DecValTok{26}\NormalTok{, }\AttributeTok{hjust =} \DecValTok{0}\NormalTok{)}
\NormalTok{    )}
\NormalTok{\}}

\CommentTok{\# Mapa facetado por fonte (num único ano)}
\NormalTok{plot\_by\_source }\OtherTok{\textless{}{-}} \ControlFlowTok{function}\NormalTok{(gdf, bra\_sf, biomes\_sf, title, }\AttributeTok{fill\_lab =} \StringTok{"kt CO2e"}\NormalTok{, }\AttributeTok{pal =}\NormalTok{ pal\_magpie, }\AttributeTok{ncol =} \DecValTok{2}\NormalTok{) \{}
  \FunctionTok{ggplot}\NormalTok{() }\SpecialCharTok{+}
    \FunctionTok{geom\_raster}\NormalTok{(}\AttributeTok{data =}\NormalTok{ gdf, }\FunctionTok{aes}\NormalTok{(lon, lat, }\AttributeTok{fill =}\NormalTok{ value)) }\SpecialCharTok{+}
    \FunctionTok{geom\_sf}\NormalTok{(}\AttributeTok{data =}\NormalTok{ biomes\_sf, }\AttributeTok{fill =} \ConstantTok{NA}\NormalTok{, }\AttributeTok{color =} \StringTok{"grey35"}\NormalTok{, }\AttributeTok{linewidth =} \FloatTok{0.25}\NormalTok{) }\SpecialCharTok{+}
    \FunctionTok{geom\_sf}\NormalTok{(}\AttributeTok{data =}\NormalTok{ bra\_sf,    }\AttributeTok{fill =} \ConstantTok{NA}\NormalTok{, }\AttributeTok{color =} \StringTok{"grey15"}\NormalTok{, }\AttributeTok{linewidth =} \FloatTok{0.4}\NormalTok{) }\SpecialCharTok{+}
    \FunctionTok{coord\_sf}\NormalTok{(}\AttributeTok{xlim =} \FunctionTok{c}\NormalTok{(}\SpecialCharTok{{-}}\DecValTok{74}\NormalTok{, }\SpecialCharTok{{-}}\DecValTok{34}\NormalTok{), }\AttributeTok{ylim =} \FunctionTok{c}\NormalTok{(}\SpecialCharTok{{-}}\DecValTok{34}\NormalTok{, }\DecValTok{6}\NormalTok{), }\AttributeTok{expand =} \ConstantTok{FALSE}\NormalTok{) }\SpecialCharTok{+}
    \FunctionTok{facet\_wrap}\NormalTok{(}\SpecialCharTok{\textasciitilde{}}\NormalTok{ source, }\AttributeTok{ncol =}\NormalTok{ ncol) }\SpecialCharTok{+}
    \FunctionTok{scale\_fill\_gradientn}\NormalTok{(}\AttributeTok{colors =}\NormalTok{ pal, }\AttributeTok{name =}\NormalTok{ fill\_lab) }\SpecialCharTok{+}
    \FunctionTok{labs}\NormalTok{(}\AttributeTok{title =}\NormalTok{ title, }\AttributeTok{caption =} \StringTok{"Source: MAgPIE / Default version (SSP2)"}\NormalTok{) }\SpecialCharTok{+}
    \FunctionTok{theme\_minimal}\NormalTok{(}\AttributeTok{base\_size =} \DecValTok{12}\NormalTok{) }\SpecialCharTok{+}
    \FunctionTok{theme}\NormalTok{(}
      \AttributeTok{legend.position =} \StringTok{"right"}\NormalTok{,}
      \AttributeTok{panel.grid      =} \FunctionTok{element\_blank}\NormalTok{(),}
      \AttributeTok{plot.title      =} \FunctionTok{element\_text}\NormalTok{(}\AttributeTok{face =} \StringTok{"bold"}\NormalTok{, }\AttributeTok{size =} \DecValTok{22}\NormalTok{)}
\NormalTok{    )}
\NormalTok{\}}
\end{Highlighting}
\end{Shaded}

\section{Introdução}\label{introduuxe7uxe3o}

A aplicação de modelos globais de uso da terra e emissões em contextos
nacionais exige cautela, dado que cada país apresenta dinâmicas próprias
em termos de atividades produtivas, estrutura fundiária, biomas e
padrões de emissões. No caso do Brasil, marcado por forte presença do
setor agropecuário e pela diversidade de ecossistemas, a incorporação
dessas especificidades é fundamental para garantir a consistência das
análises.

O MAgPIE (Model of Agricultural Production and its Impact on the
Environment) é um modelo global de equilíbrio parcial amplamente
utilizado para simular interações entre uso da terra, produção
agropecuária, comércio internacional e emissões de gases de efeito
estufa (GEE). Entretanto, sua configuração padrão (default) não
contempla ajustes específicos para a realidade brasileira, o que levanta
a necessidade de avaliar em que medida os resultados obtidos se
aproximam dos valores reportados oficialmente.

Neste relatório, busca-se avaliar a aderência do MAgPIE, em sua
configuração padrão (cenário SSP2), aos dados de emissões
disponibilizados pelo Brasil em seus inventários nacionais de GEE. A
análise concentra-se nos principais gases associados ao setor
agropecuário e ao uso da terra (metano (\(CH_4\)), óxido nitroso
(\(N_2 O\)) e dióxido de carbono (\(CO_2\))) para os anos de 2010, 2015
e 2020, conforme reportado no Quarto Inventário Nacional de Emissões e
Remoções Antrópicas, submetido à Convenção-Quadro das Nações Unidas
sobre Mudança do Clima (UNFCCC).

A comparação tem como objetivo compreender a proximidade entre as
trajetórias simuladas pelo modelo e os dados observados, destacando
potenciais divergências e padrões espaciais relevantes. A partir disso,
pretende-se identificar caminhos para ajustes metodológicos que possam
aprimorar a utilização do MAgPIE no contexto brasileiro, garantindo
maior aderência às especificidades nacionais e maior robustez nas
análises futuras.

\section{Métodos}\label{muxe9todos}

Para a obtenção dos resultados apresentados a seguir, o MAgPIE foi
executado no cenário default (SSP2), e utilizado o arquivo GDX gerado
(\texttt{fulldata.gdx}) como base para a análise. A extração no R foi
feita com as funções disponibilizadas pelo pacote \emph{magpie4}, em
especial a função \texttt{Emissions()}, que permite recuperar
diretamente do GDX as emissões de gases de efeito estufa por tipo de
gás, nível de agregação e unidade. Essa função retorna um objeto no
formato \texttt{magclass}, no qual as dimensões correspondem a regiões,
períodos de tempo e fontes de emissão. No caso deste relatório, o
parâmetro \texttt{type} foi definido de acordo com o gás de interesse
(\texttt{"ch4"}, \texttt{"n2o"}, \texttt{"co2"}), o parâmetro level foi
configurado como \texttt{"reg"} (nível regional do MAgPIE) e a unidade
escolhida foi \texttt{"GWP100AR5"}, garantindo comparabilidade com os
valores oficiais do Inventário Nacional, já que nesse caso as emissões
são reportadas em \(CO_2\) equivalente (\(CO_2e\)) conforme os fatores
de equivalência de 100 anos do AR5.

A partir dessa extração em nível regional, foi necessário redistribuir
os resultados para nível de célula, sendo 2901 células no território
brasileiro. Essa conversão foi feita com a função
\texttt{toolAggregate()} do pacote \emph{madrat}, que utiliza um mapa de
relação entre regiões e células, combinado com um peso espacial. Neste
relatório, a variável escolhida como peso foi a área de pastagem, de
forma a refletir a distribuição espacial das emissões agropecuárias.

Na dimensão das categorias de emissão, a estrutura do MAgPIE diferencia
as fontes por pares \texttt{fonte.gás} (por exemplo,
\texttt{ent\_ferm.ch4} para fermentação entérica, \texttt{awms.ch4} para
manejo de dejetos animais, ou \texttt{rice.ch4} para arroz irrigado).
Essa organização permitiu tanto a soma das emissões para obter o total
de cada gás quanto a decomposição por principais fontes. Além disso, ao
aplicar uma conversão de megatoneladas (Mt) para quilotoneladas (kt), os
resultados foram ajustados permitindo que fossem comparados diretamente
aos valores observados do Quarto Inventário Nacional, tal como reportado
pelo SIRENE.

\section{Resultados}\label{resultados}

A seguir são apresentados os resultados para metano e óxido nitroso.

\subsection{Metano}\label{metano}

A figura abaixo apresenta a distribuição espacial das emissões totais de
\(CH_4\) no Brasil (em kt \(CO_2e\), GWP100-AR5) para 2010, 2015 e 2020,
obtidas da estimação do modelo MAgPIE, em versão \emph{default}. Cada
painel corresponde a um ano; cada célula do grid resulta da
redistribuição das emissões regionais para nível de célula ponderada
pela área de pastagem. A escala de cores (tons mais escuros = valores
maiores) é comum aos três painéis, permitindo comparação visual entre
anos. O mapa mostra o total do gás (soma das fontes), sem decomposição
setorial, sobreposto aos limites de biomas.

\begin{Shaded}
\begin{Highlighting}[]
\CommentTok{\# Uso da terra (pastagem como weight)}
\NormalTok{x }\OtherTok{\textless{}{-}} \FunctionTok{read.magpie}\NormalTok{(}\StringTok{"cell.land\_0.5.mz"}\NormalTok{)[,, }\StringTok{"past"}\NormalTok{]}
\NormalTok{x[}\FunctionTok{is.na}\NormalTok{(x)] }\OtherTok{\textless{}{-}} \DecValTok{0}

\NormalTok{gas    }\OtherTok{\textless{}{-}} \StringTok{"ch4"}
\NormalTok{years  }\OtherTok{\textless{}{-}} \FunctionTok{c}\NormalTok{(}\DecValTok{2010}\NormalTok{, }\DecValTok{2015}\NormalTok{, }\DecValTok{2020}\NormalTok{)}
\NormalTok{yearS  }\OtherTok{\textless{}{-}} \FunctionTok{fmt\_year}\NormalTok{(years)                 }\CommentTok{\# "y2010.y2010", ...}
 
\CommentTok{\# 1) Emissões por gás no nível regional (reg), em Mt gás}
\NormalTok{emis\_gas      }\OtherTok{\textless{}{-}} \FunctionTok{Emissions}\NormalTok{(}\AttributeTok{gdx =} \StringTok{"fulldata.gdx"}\NormalTok{, }\AttributeTok{type =}\NormalTok{ gas, }\AttributeTok{level =} \StringTok{"reg"}\NormalTok{, }\AttributeTok{unit =} \StringTok{"GWP100AR5"}\NormalTok{, }\AttributeTok{cumulative =} \ConstantTok{FALSE}\NormalTok{)}
\NormalTok{emis\_gas }\OtherTok{\textless{}{-}}\NormalTok{ emis\_gas[,}\FunctionTok{c}\NormalTok{(}\StringTok{"y2010"}\NormalTok{, }\StringTok{"y2015"}\NormalTok{, }\StringTok{"y2020"}\NormalTok{),]}

\CommentTok{\# 2) Agrega de região {-}\textgreater{} célula usando seu mapa/weight}
\NormalTok{dis\_emis\_gas  }\OtherTok{\textless{}{-}}\NormalTok{ madrat}\SpecialCharTok{::}\FunctionTok{toolAggregate}\NormalTok{(emis\_gas, map\_file, x, }\AttributeTok{from =} \StringTok{"region"}\NormalTok{, }\AttributeTok{to =} \StringTok{"cell"}\NormalTok{)}

\CommentTok{\# 3) Para todos os \textquotesingle{}pollutants\textquotesingle{} do gás, extrai BRA x anos e empilha}
\NormalTok{ks }\OtherTok{\textless{}{-}}\NormalTok{ magclass}\SpecialCharTok{::}\FunctionTok{getItems}\NormalTok{(dis\_emis\_gas, }\AttributeTok{dim =} \DecValTok{3}\NormalTok{)}
\NormalTok{ks }\OtherTok{\textless{}{-}} \FunctionTok{grep}\NormalTok{(}\FunctionTok{paste0}\NormalTok{(}\StringTok{"}\SpecialCharTok{\textbackslash{}\textbackslash{}}\StringTok{."}\NormalTok{, gas, }\StringTok{"$"}\NormalTok{), ks, }\AttributeTok{value =} \ConstantTok{TRUE}\NormalTok{) }

\NormalTok{dfs }\OtherTok{\textless{}{-}} \FunctionTok{lapply}\NormalTok{(ks, }\ControlFlowTok{function}\NormalTok{(k) }\FunctionTok{mag\_to\_df}\NormalTok{(dis\_emis\_gas[}\StringTok{"BRA"}\NormalTok{, yearS, k]))}
\NormalTok{df\_gas\_total }\OtherTok{\textless{}{-}}\NormalTok{ dplyr}\SpecialCharTok{::}\FunctionTok{bind\_rows}\NormalTok{(dfs)}

\CommentTok{\# 4) Soma sobre as fontes para dar o total por célula/ano}
\NormalTok{df\_gas\_sum }\OtherTok{\textless{}{-}} \FunctionTok{aggregate}\NormalTok{(value }\SpecialCharTok{\textasciitilde{}}\NormalTok{ lon }\SpecialCharTok{+}\NormalTok{ lat }\SpecialCharTok{+}\NormalTok{ iso }\SpecialCharTok{+}\NormalTok{ year.start }\SpecialCharTok{+}\NormalTok{ year.end,}
                        \AttributeTok{data =}\NormalTok{ df\_gas\_total, }\AttributeTok{FUN =}\NormalTok{ sum)}

\CommentTok{\# 5) Junta com grades BRA, aplica GWP (SAR) e prepara para grid}
\NormalTok{df }\OtherTok{\textless{}{-}}\NormalTok{ BRA }\SpecialCharTok{|\textgreater{}}
  \FunctionTok{left\_join}\NormalTok{(df\_gas\_sum, }\AttributeTok{by =} \FunctionTok{c}\NormalTok{(}\StringTok{"lon"}\NormalTok{,}\StringTok{"lat"}\NormalTok{,}\StringTok{"iso"}\NormalTok{)) }\SpecialCharTok{|\textgreater{}}
  \FunctionTok{mutate}\NormalTok{(}\AttributeTok{kt\_co2eq =} \FunctionTok{mt\_to\_kt}\NormalTok{(value))}

\NormalTok{df\_map }\OtherTok{\textless{}{-}}\NormalTok{ df }\SpecialCharTok{|\textgreater{}}
\NormalTok{  dplyr}\SpecialCharTok{::}\FunctionTok{select}\NormalTok{(lon, lat, year.end, }\AttributeTok{value =}\NormalTok{ kt\_co2eq) }\SpecialCharTok{|\textgreater{}}
  \FunctionTok{filter}\NormalTok{(}\FunctionTok{is.finite}\NormalTok{(lon), }\FunctionTok{is.finite}\NormalTok{(lat), }\FunctionTok{is.finite}\NormalTok{(value))}

\CommentTok{\# 6) Rasteriza por ano e plota}
\NormalTok{g\_year }\OtherTok{\textless{}{-}} \FunctionTok{gridify\_by\_year}\NormalTok{(df\_map, }\AttributeTok{bra\_sf =}\NormalTok{ bra)}
\NormalTok{p\_ch4  }\OtherTok{\textless{}{-}} \FunctionTok{plot\_faceted\_map}\NormalTok{(g\_year, bra, biomes, }\AttributeTok{title =} \StringTok{"CH4 {-} Brasil (kt CO2e, GWP100{-}AR5)"}\NormalTok{)}
\NormalTok{p\_ch4}
\end{Highlighting}
\end{Shaded}

\begin{figure}

{\centering \includegraphics[width=1\linewidth]{emissions_files/figure-latex/ch4-1} 

}

\caption{Distribuição espacial de $CH_4$ (kt $CO_2e$, GWP100AR5) — 2010–2020.}\label{fig:ch4}
\end{figure}

A figura abaixo mostra a distribuição espacial do \(CH_4\) por fonte
para 2020 (em kt CO₂e, GWP100-AR5), destacando as duas principais
contribuições, de acordo com o MAgPIE, para o Brasil. Cada painel
corresponde a uma fonte específica entre as que mais emitem no ano
selecionado (por exemplo, fermentação entérica e manejo de
dejetos/``AWMS'', conforme o retorno do modelo), com a mesma escala de
cores nos painéis (tons mais escuros = emissões maiores).

\begin{Shaded}
\begin{Highlighting}[]
\CommentTok{\# Encontra top{-}2 (BRA, 2020)}
\NormalTok{year\_one }\OtherTok{\textless{}{-}} \FunctionTok{fmt\_year}\NormalTok{(}\DecValTok{2020}\NormalTok{)}
\NormalTok{src\_top2 }\OtherTok{\textless{}{-}} \FunctionTok{top\_k\_sources}\NormalTok{(dis\_emis\_gas, }\AttributeTok{country =} \StringTok{"BRA"}\NormalTok{, }\AttributeTok{yearS =}\NormalTok{ year\_one, }\AttributeTok{k =} \DecValTok{2}\NormalTok{,}
                          \AttributeTok{pattern =} \FunctionTok{paste0}\NormalTok{(}\StringTok{"}\SpecialCharTok{\textbackslash{}\textbackslash{}}\StringTok{."}\NormalTok{, gas, }\StringTok{"$"}\NormalTok{))}

\CommentTok{\# Extrai, aplica GWP e rasteriza por fonte}
\NormalTok{df\_src }\OtherTok{\textless{}{-}} \FunctionTok{lapply}\NormalTok{(src\_top2, }\ControlFlowTok{function}\NormalTok{(s) }\FunctionTok{extract\_one\_source}\NormalTok{(dis\_emis\_gas, }\StringTok{"BRA"}\NormalTok{, year\_one, s)) }\SpecialCharTok{|\textgreater{}}
\NormalTok{  dplyr}\SpecialCharTok{::}\FunctionTok{bind\_rows}\NormalTok{() }\SpecialCharTok{|\textgreater{}}
  \FunctionTok{mutate}\NormalTok{(}\AttributeTok{kt\_co2eq =} \FunctionTok{mt\_to\_kt}\NormalTok{(value)) }\SpecialCharTok{|\textgreater{}}
  \FunctionTok{transmute}\NormalTok{(lon, lat, source, }\AttributeTok{value =}\NormalTok{ kt\_co2eq) }\SpecialCharTok{|\textgreater{}}
  \FunctionTok{filter}\NormalTok{(}\FunctionTok{is.finite}\NormalTok{(lon), }\FunctionTok{is.finite}\NormalTok{(lat), }\FunctionTok{is.finite}\NormalTok{(value))}

\NormalTok{g\_src }\OtherTok{\textless{}{-}} \FunctionTok{split}\NormalTok{(df\_src, df\_src}\SpecialCharTok{$}\NormalTok{source) }\SpecialCharTok{|\textgreater{}}
  \FunctionTok{lapply}\NormalTok{(}\ControlFlowTok{function}\NormalTok{(d) \{}
\NormalTok{    r  }\OtherTok{\textless{}{-}}\NormalTok{ terra}\SpecialCharTok{::}\FunctionTok{rast}\NormalTok{(d[, }\FunctionTok{c}\NormalTok{(}\StringTok{"lon"}\NormalTok{,}\StringTok{"lat"}\NormalTok{,}\StringTok{"value"}\NormalTok{)], }\AttributeTok{type =} \StringTok{"xyz"}\NormalTok{, }\AttributeTok{crs =} \StringTok{"EPSG:4326"}\NormalTok{)}
\NormalTok{    r  }\OtherTok{\textless{}{-}}\NormalTok{ terra}\SpecialCharTok{::}\FunctionTok{crop}\NormalTok{(r, terra}\SpecialCharTok{::}\FunctionTok{vect}\NormalTok{(bra))}
\NormalTok{    r  }\OtherTok{\textless{}{-}}\NormalTok{ terra}\SpecialCharTok{::}\FunctionTok{mask}\NormalTok{(r, terra}\SpecialCharTok{::}\FunctionTok{vect}\NormalTok{(bra))}
\NormalTok{    out }\OtherTok{\textless{}{-}} \FunctionTok{as.data.frame}\NormalTok{(r, }\AttributeTok{xy =} \ConstantTok{TRUE}\NormalTok{, }\AttributeTok{na.rm =} \ConstantTok{TRUE}\NormalTok{)}
    \FunctionTok{names}\NormalTok{(out) }\OtherTok{\textless{}{-}} \FunctionTok{c}\NormalTok{(}\StringTok{"lon"}\NormalTok{,}\StringTok{"lat"}\NormalTok{,}\StringTok{"value"}\NormalTok{)}
\NormalTok{    out}\SpecialCharTok{$}\NormalTok{source }\OtherTok{\textless{}{-}} \FunctionTok{unique}\NormalTok{(d}\SpecialCharTok{$}\NormalTok{source)}
\NormalTok{    out}
\NormalTok{  \}) }\SpecialCharTok{|\textgreater{}}
\NormalTok{  dplyr}\SpecialCharTok{::}\FunctionTok{bind\_rows}\NormalTok{()}

\NormalTok{labels\_dict }\OtherTok{\textless{}{-}} \FunctionTok{c}\NormalTok{(}
  \StringTok{"awms.ch4"}     \OtherTok{=} \StringTok{"Manejo de dejetos animais"}\NormalTok{,}
  \StringTok{"ent\_ferm.ch4"} \OtherTok{=} \StringTok{"Fermentação entérica"}\NormalTok{,}
  \StringTok{"rice.ch4"}     \OtherTok{=} \StringTok{"Arroz irrigado"}
\NormalTok{)}
\NormalTok{lab }\OtherTok{\textless{}{-}}\NormalTok{ labels\_dict[src\_top2]; lab[}\FunctionTok{is.na}\NormalTok{(lab)] }\OtherTok{\textless{}{-}}\NormalTok{ src\_top2[}\FunctionTok{is.na}\NormalTok{(lab)]}
\NormalTok{g\_src}\SpecialCharTok{$}\NormalTok{source }\OtherTok{\textless{}{-}} \FunctionTok{factor}\NormalTok{(g\_src}\SpecialCharTok{$}\NormalTok{source, }\AttributeTok{levels =}\NormalTok{ src\_top2, }\AttributeTok{labels =}\NormalTok{ lab)}

\NormalTok{p\_ch4\_src }\OtherTok{\textless{}{-}} \FunctionTok{plot\_by\_source}\NormalTok{(g\_src, bra, biomes, }\AttributeTok{title =} \StringTok{"CH4 {-} Brasil (kt CO2e, GWP100{-}AR5) — principais fontes (2020)"}\NormalTok{)}
\NormalTok{p\_ch4\_src}
\end{Highlighting}
\end{Shaded}

\begin{figure}

{\centering \includegraphics[width=1\linewidth]{emissions_files/figure-latex/source-ch4-1} 

}

\caption{Distribuição espacial de $CH_4$ considerando as principais fontes de emissão (kt $CO_2e$, GWP100AR5) — 2020.}\label{fig:source-ch4}
\end{figure}

A figura apresenta a comparação entre os totais de emissões de \(CH_4\)
no Brasil, conforme estimados pelo modelo MAgPIE em sua versão default
(SSP2), e os valores oficiais reportados no Quarto Inventário Nacional
de Emissões, para os anos de 2010, 2015 e 2020.

\begin{Shaded}
\begin{Highlighting}[]
\CommentTok{\# Totais CH4 Brasil a partir do MAgPIE (kt CO2e, AR5)}
\NormalTok{magpie\_br }\OtherTok{\textless{}{-}}\NormalTok{ df\_map }\SpecialCharTok{|\textgreater{}}
\NormalTok{  dplyr}\SpecialCharTok{::}\FunctionTok{mutate}\NormalTok{(}\AttributeTok{year =} \FunctionTok{as.integer}\NormalTok{(}\FunctionTok{gsub}\NormalTok{(}\StringTok{"[\^{}0{-}9]"}\NormalTok{, }\StringTok{""}\NormalTok{, year.end))) }\SpecialCharTok{|\textgreater{}}
\NormalTok{  dplyr}\SpecialCharTok{::}\FunctionTok{filter}\NormalTok{(year }\SpecialCharTok{\%in\%} \FunctionTok{c}\NormalTok{(}\DecValTok{2010}\NormalTok{, }\DecValTok{2015}\NormalTok{, }\DecValTok{2020}\NormalTok{)) }\SpecialCharTok{|\textgreater{}}
\NormalTok{  dplyr}\SpecialCharTok{::}\FunctionTok{group\_by}\NormalTok{(year) }\SpecialCharTok{|\textgreater{}}
\NormalTok{  dplyr}\SpecialCharTok{::}\FunctionTok{summarise}\NormalTok{(}\AttributeTok{value\_kt =} \FunctionTok{sum}\NormalTok{(value, }\AttributeTok{na.rm =} \ConstantTok{TRUE}\NormalTok{), }\AttributeTok{.groups =} \StringTok{"drop"}\NormalTok{) }\SpecialCharTok{|\textgreater{}}
\NormalTok{  dplyr}\SpecialCharTok{::}\FunctionTok{mutate}\NormalTok{(}\AttributeTok{source =} \StringTok{"MAgPIE"}\NormalTok{)}

\CommentTok{\# Totais CH4 Brasil do SIRENE}
\NormalTok{sirene\_br }\OtherTok{\textless{}{-}}\NormalTok{ tibble}\SpecialCharTok{::}\FunctionTok{tibble}\NormalTok{(}
  \AttributeTok{year =} \FunctionTok{c}\NormalTok{(}\DecValTok{2010}\NormalTok{, }\DecValTok{2015}\NormalTok{, }\DecValTok{2020}\NormalTok{),}
  \AttributeTok{value\_kt =} \FunctionTok{c}\NormalTok{(}\DecValTok{432565}\NormalTok{, }\DecValTok{434079}\NormalTok{, }\DecValTok{449829}\NormalTok{),}
  \AttributeTok{source =} \StringTok{"SIRENE"}
\NormalTok{)}

\CommentTok{\# Junta e plota}
\NormalTok{df\_comp }\OtherTok{\textless{}{-}}\NormalTok{ dplyr}\SpecialCharTok{::}\FunctionTok{bind\_rows}\NormalTok{(magpie\_br, sirene\_br)}

\FunctionTok{ggplot}\NormalTok{(df\_comp, }\FunctionTok{aes}\NormalTok{(}\AttributeTok{x =} \FunctionTok{factor}\NormalTok{(year), }\AttributeTok{y =}\NormalTok{ value\_kt, }\AttributeTok{fill =}\NormalTok{ source)) }\SpecialCharTok{+}
  \FunctionTok{geom\_col}\NormalTok{(}\AttributeTok{position =} \FunctionTok{position\_dodge}\NormalTok{(}\AttributeTok{width =} \FloatTok{0.8}\NormalTok{), }\AttributeTok{width =} \FloatTok{0.75}\NormalTok{) }\SpecialCharTok{+}
  \FunctionTok{geom\_text}\NormalTok{(}\FunctionTok{aes}\NormalTok{(}\AttributeTok{label =}\NormalTok{ scales}\SpecialCharTok{::}\FunctionTok{comma}\NormalTok{(value\_kt)),}
            \AttributeTok{position =} \FunctionTok{position\_dodge}\NormalTok{(}\AttributeTok{width =} \FloatTok{0.8}\NormalTok{),}
            \AttributeTok{vjust =} \SpecialCharTok{{-}}\FloatTok{0.2}\NormalTok{, }\AttributeTok{size =} \FloatTok{3.5}\NormalTok{) }\SpecialCharTok{+}
  \FunctionTok{scale\_y\_continuous}\NormalTok{(}\AttributeTok{labels =}\NormalTok{ scales}\SpecialCharTok{::}\FunctionTok{label\_comma}\NormalTok{()) }\SpecialCharTok{+}
  \FunctionTok{scale\_fill\_manual}\NormalTok{(}\AttributeTok{values =} \FunctionTok{c}\NormalTok{(}\StringTok{"MAgPIE"} \OtherTok{=}\NormalTok{ pal\_magpie[}\DecValTok{5}\NormalTok{], }\StringTok{"SIRENE"} \OtherTok{=}\NormalTok{ pal\_magpie[}\DecValTok{15}\NormalTok{])) }\SpecialCharTok{+}
  \FunctionTok{labs}\NormalTok{(}
    \AttributeTok{title =} \StringTok{"CH4 — Brasil (kt CO2e, GWP100{-}AR5): MAgPIE vs. SIRENE)"}\NormalTok{,}
    \AttributeTok{x =} \ConstantTok{NULL}\NormalTok{, }\AttributeTok{y =} \StringTok{"kt CO2e"}\NormalTok{,}
\NormalTok{  ) }\SpecialCharTok{+}
  \FunctionTok{theme\_minimal}\NormalTok{(}\AttributeTok{base\_size =} \DecValTok{12}\NormalTok{) }\SpecialCharTok{+}
  \FunctionTok{theme}\NormalTok{(}\AttributeTok{legend.position =} \StringTok{"top"}\NormalTok{)}
\end{Highlighting}
\end{Shaded}

\begin{figure}

{\centering \includegraphics[width=1\linewidth]{emissions_files/figure-latex/totais-ch4-1} 

}

\caption{Totais de emissões de $CH_4$ no Brasil (kt $CO_2e$, GWP100-AR5) em 2010, 2015 e 2020, comparando os resultados do modelo MAgPIE (cenário default, SSP2) com os valores reportados pelo Quarto Inventário Nacional de Emissões (SIRENE).}\label{fig:totais-ch4}
\end{figure}

\subsection{Óxido Nitroso}\label{uxf3xido-nitroso}

A figura abaixo apresenta a distribuição espacial das emissões totais de
\(N_2O\) no Brasil (em kt CO₂e, GWP100--AR5), conforme estimativas do
modelo MAgPIE em sua configuração padrão (default) para os anos de 2010,
2015 e 2020. Diferentemente do caso das emissões de metano, cada célula
do grid resulta da redistribuição das emissões regionais para o nível de
célula, ponderada pela área de culturas agrícolas.

De acordo com as estimativas do MAgPIE, a maior concentração das
emissões de \(N_2O\) está localizada principalmente na região
Centro-Sul, abrangendo áreas dos estados de Mato Grosso do Sul, Paraná,
São Paulo, Minas Gerais e parte de Goiás.

\begin{Shaded}
\begin{Highlighting}[]
\CommentTok{\# Uso da terra (pastagem como weight)}
\NormalTok{x }\OtherTok{\textless{}{-}} \FunctionTok{read.magpie}\NormalTok{(}\StringTok{"cell.land\_0.5.mz"}\NormalTok{)[,, }\StringTok{"crop"}\NormalTok{]}
\NormalTok{x[}\FunctionTok{is.na}\NormalTok{(x)] }\OtherTok{\textless{}{-}} \DecValTok{0}

\NormalTok{pal\_magpie }\OtherTok{\textless{}{-}} \FunctionTok{hcl.colors}\NormalTok{(}\DecValTok{20}\NormalTok{, }\StringTok{"Reds"}\NormalTok{, }\AttributeTok{rev =} \ConstantTok{TRUE}\NormalTok{)}

\NormalTok{gas    }\OtherTok{\textless{}{-}} \StringTok{"n2o"}
\NormalTok{years  }\OtherTok{\textless{}{-}} \FunctionTok{c}\NormalTok{(}\DecValTok{2010}\NormalTok{, }\DecValTok{2015}\NormalTok{, }\DecValTok{2020}\NormalTok{)}
\NormalTok{yearS  }\OtherTok{\textless{}{-}} \FunctionTok{fmt\_year}\NormalTok{(years)                 }\CommentTok{\# "y2010.y2010", ...}

\CommentTok{\# 1) Emissões por gás no nível regional (reg), em Mt gás}
\NormalTok{emis\_gas      }\OtherTok{\textless{}{-}} \FunctionTok{Emissions}\NormalTok{(}\AttributeTok{gdx =} \StringTok{"fulldata.gdx"}\NormalTok{, }\AttributeTok{type =}\NormalTok{ gas, }\AttributeTok{level =} \StringTok{"reg"}\NormalTok{, }\AttributeTok{unit =} \StringTok{"GWP100AR5"}\NormalTok{)}
\NormalTok{emis\_gas }\OtherTok{\textless{}{-}}\NormalTok{ emis\_gas[,}\FunctionTok{c}\NormalTok{(}\StringTok{"y2010"}\NormalTok{, }\StringTok{"y2015"}\NormalTok{, }\StringTok{"y2020"}\NormalTok{),]}

\CommentTok{\# 2) Agrega de região {-}\textgreater{} célula usando seu mapa/weight}
\NormalTok{dis\_emis\_gas  }\OtherTok{\textless{}{-}}\NormalTok{ madrat}\SpecialCharTok{::}\FunctionTok{toolAggregate}\NormalTok{(emis\_gas, map\_file, x, }\AttributeTok{from =} \StringTok{"region"}\NormalTok{, }\AttributeTok{to =} \StringTok{"cell"}\NormalTok{)}

\CommentTok{\# 3) Para todos os \textquotesingle{}pollutants\textquotesingle{} do gás, extrai BRA x anos e empilha}
\NormalTok{ks }\OtherTok{\textless{}{-}}\NormalTok{ magclass}\SpecialCharTok{::}\FunctionTok{getItems}\NormalTok{(dis\_emis\_gas, }\AttributeTok{dim =} \DecValTok{3}\NormalTok{)}
\NormalTok{ks }\OtherTok{\textless{}{-}} \FunctionTok{grep}\NormalTok{(}\FunctionTok{paste0}\NormalTok{(}\StringTok{"}\SpecialCharTok{\textbackslash{}\textbackslash{}}\StringTok{."}\NormalTok{, gas, }\StringTok{"$"}\NormalTok{), ks, }\AttributeTok{value =} \ConstantTok{TRUE}\NormalTok{) }

\NormalTok{dfs\_n }\OtherTok{\textless{}{-}} \FunctionTok{lapply}\NormalTok{(ks, }\ControlFlowTok{function}\NormalTok{(k) }\FunctionTok{mag\_to\_df}\NormalTok{(dis\_emis\_gas[}\StringTok{"BRA"}\NormalTok{, yearS, k]))}
\NormalTok{df\_gas\_total }\OtherTok{\textless{}{-}}\NormalTok{ dplyr}\SpecialCharTok{::}\FunctionTok{bind\_rows}\NormalTok{(dfs\_n)}

\CommentTok{\# 4) Soma sobre as fontes para dar o total por célula/ano}
\NormalTok{df\_gas\_sum }\OtherTok{\textless{}{-}} \FunctionTok{aggregate}\NormalTok{(value }\SpecialCharTok{\textasciitilde{}}\NormalTok{ lon }\SpecialCharTok{+}\NormalTok{ lat }\SpecialCharTok{+}\NormalTok{ iso }\SpecialCharTok{+}\NormalTok{ year.start }\SpecialCharTok{+}\NormalTok{ year.end,}
                        \AttributeTok{data =}\NormalTok{ df\_gas\_total, }\AttributeTok{FUN =}\NormalTok{ sum)}

\CommentTok{\# 5) Junta com grades BRA, aplica GWP (SAR) e prepara para grid}
\NormalTok{df\_n2o }\OtherTok{\textless{}{-}}\NormalTok{ BRA }\SpecialCharTok{|\textgreater{}}
  \FunctionTok{left\_join}\NormalTok{(df\_gas\_sum, }\AttributeTok{by =} \FunctionTok{c}\NormalTok{(}\StringTok{"lon"}\NormalTok{,}\StringTok{"lat"}\NormalTok{,}\StringTok{"iso"}\NormalTok{)) }\SpecialCharTok{|\textgreater{}}
  \FunctionTok{mutate}\NormalTok{(}\AttributeTok{kt\_co2eq =} \FunctionTok{mt\_to\_kt}\NormalTok{(value))}

\NormalTok{df\_map\_n2o }\OtherTok{\textless{}{-}}\NormalTok{ df\_n2o }\SpecialCharTok{|\textgreater{}}
\NormalTok{  dplyr}\SpecialCharTok{::}\FunctionTok{select}\NormalTok{(lon, lat, year.end, }\AttributeTok{value =}\NormalTok{ kt\_co2eq) }\SpecialCharTok{|\textgreater{}}
  \FunctionTok{filter}\NormalTok{(}\FunctionTok{is.finite}\NormalTok{(lon), }\FunctionTok{is.finite}\NormalTok{(lat), }\FunctionTok{is.finite}\NormalTok{(value))}

\CommentTok{\# 6) Rasteriza por ano e plota}
\NormalTok{g\_year\_n2o }\OtherTok{\textless{}{-}} \FunctionTok{gridify\_by\_year}\NormalTok{(df\_map\_n2o, }\AttributeTok{bra\_sf =}\NormalTok{ bra)}
\NormalTok{p\_n2o  }\OtherTok{\textless{}{-}} \FunctionTok{plot\_faceted\_map}\NormalTok{(g\_year\_n2o, bra, biomes, }\AttributeTok{title =} \StringTok{"N2O {-} Brasil (kt CO2e, GWP100{-}AR5"}\NormalTok{)}
\NormalTok{p\_n2o}
\end{Highlighting}
\end{Shaded}

\begin{figure}

{\centering \includegraphics[width=1\linewidth]{emissions_files/figure-latex/unnamed-chunk-1-1} 

}

\caption{Distribuição espacial de $N_2O$ (kt $CO_2e$, GWP100AR5) — 2010–2020.}\label{fig:unnamed-chunk-1}
\end{figure}

O modelo incorpora múltiplas fontes de \(N_2O\) (33 fontes), incluindo:
matéria orgânica do solo, fertilizantes inorgânicos e esterco aplicados
em lavouras e pastagens, manejo de resíduos animais, decomposição e
queima de resíduos agrícolas, cultivo de arroz, além de óxido nitroso
proveniente de carbono vegetal, serapilheira e outras vegetações,
\emph{etc}.

A figura a seguir apresenta a distribuição espacial do \(N_2O\) no
Brasil em 2020 (kt CO2e, GWP100-AR5), considerando as duas principais
fontes de emissão segundo o modelo MAgPIE: \textbf{esterco em pastagens}
e \textbf{manejo de dejetos animais}. De acordo com as estimativas, em
2020, o Brasil teria emitido 113.298,6 ktCO\_2e de óxido nitroso. Deste
total, 26\% são atribuídos ao esterco em pastagens e 20\% ao manejo de
dejetos animais.

As bases de dados utilizadas incluem FAO, Associação Internacional de
Fertilizantes (IFA) e literatura científica especializada, com detalhes
metodológicos descritos em Bodirsky et al.~(2012).

\begin{Shaded}
\begin{Highlighting}[]
\CommentTok{\# Encontra top{-}2 n2o (BRA, 2020)}
\NormalTok{year\_one\_n }\OtherTok{\textless{}{-}} \FunctionTok{fmt\_year}\NormalTok{(}\DecValTok{2020}\NormalTok{)}
\NormalTok{src\_top2\_n }\OtherTok{\textless{}{-}} \FunctionTok{top\_k\_sources}\NormalTok{(dis\_emis\_gas, }\AttributeTok{country =} \StringTok{"BRA"}\NormalTok{, }\AttributeTok{yearS =}\NormalTok{ year\_one\_n, }\AttributeTok{k =} \DecValTok{2}\NormalTok{,}
                          \AttributeTok{pattern =} \FunctionTok{paste0}\NormalTok{(}\StringTok{"}\SpecialCharTok{\textbackslash{}\textbackslash{}}\StringTok{."}\NormalTok{, gas, }\StringTok{"$"}\NormalTok{))}

\CommentTok{\# Extrai, aplica GWP e rasteriza por fonte}
\NormalTok{df\_src\_n }\OtherTok{\textless{}{-}} \FunctionTok{lapply}\NormalTok{(src\_top2\_n, }\ControlFlowTok{function}\NormalTok{(s) }\FunctionTok{extract\_one\_source}\NormalTok{(dis\_emis\_gas, }\StringTok{"BRA"}\NormalTok{, }
\NormalTok{                                                              year\_one\_n, s)) }\SpecialCharTok{|\textgreater{}}
\NormalTok{  dplyr}\SpecialCharTok{::}\FunctionTok{bind\_rows}\NormalTok{() }\SpecialCharTok{|\textgreater{}}
  \FunctionTok{mutate}\NormalTok{(}\AttributeTok{kt\_co2eq =} \FunctionTok{mt\_to\_kt}\NormalTok{(value)) }\SpecialCharTok{|\textgreater{}}
  \FunctionTok{transmute}\NormalTok{(lon, lat, source, }\AttributeTok{value =}\NormalTok{ kt\_co2eq) }\SpecialCharTok{|\textgreater{}}
  \FunctionTok{filter}\NormalTok{(}\FunctionTok{is.finite}\NormalTok{(lon), }\FunctionTok{is.finite}\NormalTok{(lat), }\FunctionTok{is.finite}\NormalTok{(value))}

\NormalTok{g\_src\_n }\OtherTok{\textless{}{-}} \FunctionTok{split}\NormalTok{(df\_src\_n, df\_src\_n}\SpecialCharTok{$}\NormalTok{source) }\SpecialCharTok{|\textgreater{}}
  \FunctionTok{lapply}\NormalTok{(}\ControlFlowTok{function}\NormalTok{(d) \{}
\NormalTok{    r  }\OtherTok{\textless{}{-}}\NormalTok{ terra}\SpecialCharTok{::}\FunctionTok{rast}\NormalTok{(d[, }\FunctionTok{c}\NormalTok{(}\StringTok{"lon"}\NormalTok{,}\StringTok{"lat"}\NormalTok{,}\StringTok{"value"}\NormalTok{)], }\AttributeTok{type =} \StringTok{"xyz"}\NormalTok{, }\AttributeTok{crs =} \StringTok{"EPSG:4326"}\NormalTok{)}
\NormalTok{    r  }\OtherTok{\textless{}{-}}\NormalTok{ terra}\SpecialCharTok{::}\FunctionTok{crop}\NormalTok{(r, terra}\SpecialCharTok{::}\FunctionTok{vect}\NormalTok{(bra))}
\NormalTok{    r  }\OtherTok{\textless{}{-}}\NormalTok{ terra}\SpecialCharTok{::}\FunctionTok{mask}\NormalTok{(r, terra}\SpecialCharTok{::}\FunctionTok{vect}\NormalTok{(bra))}
\NormalTok{    out }\OtherTok{\textless{}{-}} \FunctionTok{as.data.frame}\NormalTok{(r, }\AttributeTok{xy =} \ConstantTok{TRUE}\NormalTok{, }\AttributeTok{na.rm =} \ConstantTok{TRUE}\NormalTok{)}
    \FunctionTok{names}\NormalTok{(out) }\OtherTok{\textless{}{-}} \FunctionTok{c}\NormalTok{(}\StringTok{"lon"}\NormalTok{,}\StringTok{"lat"}\NormalTok{,}\StringTok{"value"}\NormalTok{)}
\NormalTok{    out}\SpecialCharTok{$}\NormalTok{source }\OtherTok{\textless{}{-}} \FunctionTok{unique}\NormalTok{(d}\SpecialCharTok{$}\NormalTok{source)}
\NormalTok{    out}
\NormalTok{  \}) }\SpecialCharTok{|\textgreater{}}
\NormalTok{  dplyr}\SpecialCharTok{::}\FunctionTok{bind\_rows}\NormalTok{()}

\NormalTok{labels\_dict }\OtherTok{\textless{}{-}} \FunctionTok{c}\NormalTok{(}
  \StringTok{"man\_past.n2o"} \OtherTok{=} \StringTok{"Esterco em pastagens"}\NormalTok{,}
  \StringTok{"awms.n2o"}  \OtherTok{=} \StringTok{"Manejo de dejetos animais"}
\NormalTok{  )}

\NormalTok{lab }\OtherTok{\textless{}{-}}\NormalTok{ labels\_dict[src\_top2\_n]; lab[}\FunctionTok{is.na}\NormalTok{(lab)] }\OtherTok{\textless{}{-}}\NormalTok{ src\_top2\_n[}\FunctionTok{is.na}\NormalTok{(lab)]}
\NormalTok{g\_src\_n}\SpecialCharTok{$}\NormalTok{source }\OtherTok{\textless{}{-}} \FunctionTok{factor}\NormalTok{(g\_src\_n}\SpecialCharTok{$}\NormalTok{source, }\AttributeTok{levels =}\NormalTok{ src\_top2\_n, }\AttributeTok{labels =}\NormalTok{ lab)}

\NormalTok{p\_n2o\_src }\OtherTok{\textless{}{-}} \FunctionTok{plot\_by\_source}\NormalTok{(g\_src\_n, bra, biomes, }\AttributeTok{title =} \StringTok{"N2O {-} Brasil (kt CO2e, GWP100{-}AR5) {-} principais fontes (2020)"}\NormalTok{)}
\NormalTok{p\_n2o\_src}
\end{Highlighting}
\end{Shaded}

\begin{figure}

{\centering \includegraphics[width=1\linewidth]{emissions_files/figure-latex/unnamed-chunk-2-1} 

}

\caption{Distribuição espacial de $N_2O$ considerando as principais fontes de emissão (kt $CO_2e$, GWP100AR5) — 2020.}\label{fig:unnamed-chunk-2}
\end{figure}

A figura compara os totais de emissão de \(N_2O\) no Brasil estimados
pelo modelo MAgPIE (versão default, \(SSP2\)) com os valores oficiais do
Quarto Inventário Nacional de Emissões para os anos de 2010, 2015 e
2020. De modo geral, as estimativas ficaram relativamente próximas.
Contudo, é importante considerar que a desagregação das atividades
relacionadas às emissões são feitas de formas distintas, o que pode
levar a conclusões também distintas quanto às principais fontes de
emissões.

No SIRENE, dentro do setor agropecuário, a principal fonte de emissão é
o manejo de solos agrícolas, que inclui fertilizantes sintéticos, adubos
orgânicos, deposição direta de dejetos, resíduos agrícolas,
mineralização de nitrogênio associada à perda de carbono do solo e
manejo de solos orgânicos. O esterco em pastagens, identificado pelo
MAgPIE como uma das principais fontes, está incluído nesse agrupamento.

As bases de dados e referências utilizadas pelo MAgPIE e pelo SIRENE
diferem significativamente, embora ambos adotem fatores de emissão
conforme as diretrizes do IPCC (2006).

\begin{Shaded}
\begin{Highlighting}[]
\CommentTok{\# Totais N20 Brasil a partir do MAgPIE (kt CO2e, AR5)}
\NormalTok{magpie\_br\_n }\OtherTok{\textless{}{-}}\NormalTok{ df\_map\_n2o }\SpecialCharTok{|\textgreater{}}
\NormalTok{  dplyr}\SpecialCharTok{::}\FunctionTok{mutate}\NormalTok{(}\AttributeTok{year =} \FunctionTok{as.integer}\NormalTok{(}\FunctionTok{gsub}\NormalTok{(}\StringTok{"[\^{}0{-}9]"}\NormalTok{, }\StringTok{""}\NormalTok{, year.end))) }\SpecialCharTok{|\textgreater{}}
\NormalTok{  dplyr}\SpecialCharTok{::}\FunctionTok{filter}\NormalTok{(year }\SpecialCharTok{\%in\%} \FunctionTok{c}\NormalTok{(}\DecValTok{2010}\NormalTok{, }\DecValTok{2015}\NormalTok{, }\DecValTok{2020}\NormalTok{)) }\SpecialCharTok{|\textgreater{}}
\NormalTok{  dplyr}\SpecialCharTok{::}\FunctionTok{group\_by}\NormalTok{(year) }\SpecialCharTok{|\textgreater{}}
\NormalTok{  dplyr}\SpecialCharTok{::}\FunctionTok{summarise}\NormalTok{(}\AttributeTok{value\_kt =} \FunctionTok{sum}\NormalTok{(value, }\AttributeTok{na.rm =} \ConstantTok{TRUE}\NormalTok{), }\AttributeTok{.groups =} \StringTok{"drop"}\NormalTok{) }\SpecialCharTok{|\textgreater{}}
\NormalTok{  dplyr}\SpecialCharTok{::}\FunctionTok{mutate}\NormalTok{(}\AttributeTok{source =} \StringTok{"MAgPIE"}\NormalTok{)}

\CommentTok{\# Totais N2O Brasil do SIRENE}
\NormalTok{sirene\_br\_n }\OtherTok{\textless{}{-}}\NormalTok{ tibble}\SpecialCharTok{::}\FunctionTok{tibble}\NormalTok{(}
  \AttributeTok{year =} \FunctionTok{c}\NormalTok{(}\DecValTok{2010}\NormalTok{, }\DecValTok{2015}\NormalTok{, }\DecValTok{2020}\NormalTok{),}
  \AttributeTok{value\_kt =} \FunctionTok{c}\NormalTok{(}\DecValTok{130140}\NormalTok{,  }\DecValTok{140231}\NormalTok{,  }\DecValTok{156306}\NormalTok{), }\CommentTok{\# excluir, energia, Resíduos, processos industriais}
  \AttributeTok{source =} \StringTok{"SIRENE"}
\NormalTok{  )}

\CommentTok{\# Junta e plota}
\NormalTok{df\_comp\_n }\OtherTok{\textless{}{-}}\NormalTok{ dplyr}\SpecialCharTok{::}\FunctionTok{bind\_rows}\NormalTok{(magpie\_br\_n, sirene\_br\_n)}

\FunctionTok{ggplot}\NormalTok{(df\_comp\_n, }\FunctionTok{aes}\NormalTok{(}\AttributeTok{x =} \FunctionTok{factor}\NormalTok{(year), }\AttributeTok{y =}\NormalTok{ value\_kt, }\AttributeTok{fill =}\NormalTok{ source)) }\SpecialCharTok{+}
  \FunctionTok{geom\_col}\NormalTok{(}\AttributeTok{position =} \FunctionTok{position\_dodge}\NormalTok{(}\AttributeTok{width =} \FloatTok{0.8}\NormalTok{), }\AttributeTok{width =} \FloatTok{0.75}\NormalTok{) }\SpecialCharTok{+}
  \FunctionTok{geom\_text}\NormalTok{(}\FunctionTok{aes}\NormalTok{(}\AttributeTok{label =}\NormalTok{ scales}\SpecialCharTok{::}\FunctionTok{comma}\NormalTok{(value\_kt)),}
            \AttributeTok{position =} \FunctionTok{position\_dodge}\NormalTok{(}\AttributeTok{width =} \FloatTok{0.8}\NormalTok{),}
            \AttributeTok{vjust =} \SpecialCharTok{{-}}\FloatTok{0.2}\NormalTok{, }\AttributeTok{size =} \FloatTok{3.5}\NormalTok{) }\SpecialCharTok{+}
  \FunctionTok{scale\_y\_continuous}\NormalTok{(}\AttributeTok{labels =}\NormalTok{ scales}\SpecialCharTok{::}\FunctionTok{label\_comma}\NormalTok{()) }\SpecialCharTok{+}
  \FunctionTok{scale\_fill\_manual}\NormalTok{(}\AttributeTok{values =} \FunctionTok{c}\NormalTok{(}\StringTok{"MAgPIE"} \OtherTok{=}\NormalTok{ pal\_magpie[}\DecValTok{5}\NormalTok{], }\StringTok{"SIRENE"} \OtherTok{=}\NormalTok{ pal\_magpie[}\DecValTok{15}\NormalTok{])) }\SpecialCharTok{+}
  \FunctionTok{labs}\NormalTok{(}
    \AttributeTok{title =} \StringTok{"N20 {-} Brasil (kt CO2e, GWP100{-}AR5): MAgPIE vs. SIRENE"}\NormalTok{,}
    \AttributeTok{x =} \ConstantTok{NULL}\NormalTok{, }\AttributeTok{y =} \StringTok{"kt CO2e"}\NormalTok{,}
\NormalTok{  ) }\SpecialCharTok{+}
  \FunctionTok{theme\_minimal}\NormalTok{(}\AttributeTok{base\_size =} \DecValTok{12}\NormalTok{) }\SpecialCharTok{+}
  \FunctionTok{theme}\NormalTok{(}\AttributeTok{legend.position =} \StringTok{"top"}\NormalTok{)}
\end{Highlighting}
\end{Shaded}

\begin{figure}

{\centering \includegraphics[width=1\linewidth]{emissions_files/figure-latex/unnamed-chunk-3-1} 

}

\caption{Totais de emissões de $N_2O$ no Brasil (kt $CO_2e$, GWP100-AR5) em 2010, 2015 e 2020, comparando os resultados do modelo MAgPIE (cenário default, SSP2) com os valores reportados pelo Quarto Inventário Nacional de Emissões (SIRENE).}\label{fig:unnamed-chunk-3}
\end{figure}

\end{document}
